\section{The recognition hypotheses}
\label{sec:recognition-sources}
% Each carrier is a typed structure in Lean whose arithmetic content is a
% hypothesis field (lean_coverage = recognition_source).

In the Six Birds framework a \emph{recognition source} is arithmetic
content that a construction rests on but does not derive: it is stated
precisely, attached to the object, and used as a hypothesis.  In this
paper there are three, one for each of three factors of the BSD formula
that the formal construction does not supply by itself.
\begin{itemize}
\item $\GammaPD$ concerns the local Tamagawa factor at an odd prime of
  additive reduction;
\item $\GammaSP$ concerns the $\Sha$ factor, through a signed local
  comparison;
\item $\GammaGZ$ concerns the leading term in analytic rank at least $2$.
\end{itemize}
\Cref{tab:recognition-sources} summarizes them.  We say at once what will
matter for the main theorem: its scalar conclusion uses only $\GammaGZ$ in
rank at least $2$, and the classical identity in rank at most $1$
(\Cref{sec:shell-readout}).  The two local hypotheses are part of the
recognition predicate, with their scopes recorded below, but the scalar
argument does not use them.
\begin{table}[tbp]
\centering
\footnotesize
\caption{The three recognition hypotheses.  The last column records
whether the scalar argument of \Cref{thm:closure:strong-bsd-conditional}
uses the hypothesis.}
\label{tab:recognition-sources}
\begin{tabularx}{\textwidth}{@{}>{\raggedright\arraybackslash}p{0.17\textwidth}
>{\raggedright\arraybackslash}p{0.25\textwidth}
>{\raggedright\arraybackslash}X
>{\raggedright\arraybackslash}p{0.12\textwidth}@{}}
\toprule
Hypothesis & Scope & Content & Used for scalar BSD \\
\midrule
\(\GammaPD\) & odd prime \(p\), additive type II or III, \(c_p\in\{1,2\}\),
\(\Sha[p^\infty]=0\), tame \(L/\Qp\) giving good reduction & descent congruence
modulo \(\Zp^\times\) with factor \(c_p^{-1}\) and a unit; a local factor
matched with \(\Tam(E)\) & no \\
\addlinespace
\(\GammaSP\) & prime \(p\) of additive reduction in the signed scope, tame \(L/\Qp\) giving good reduction & signed corestriction determinant \(\equiv c_p\)
modulo \(\Zp^\times\); a factor matched with \(\#\Sha(E/\Q)\) & no \\
\addlinespace
\(\GammaGZ\) & no CM, analytic rank \(r\ge2\) & \(\psi_-(E,r)=0\), with
\(\kapr=\#\Sha/\#E(\Q)_{\tors}^{\,2}\) and matched factors; equivalent to
Strong BSD for \(E\) when \(\Sha\) is finite & yes, in rank \(\ge2\) \\
\bottomrule
\end{tabularx}
\end{table}


\begin{definition}[$\GammaPD$: $p$-adic descent hypothesis]\label{def:closure:gamma-padic-descent}
Let $E/\Q$ and let $p$ be an odd prime at which $E$ has additive
reduction of Kodaira type II or III, with $c_p(E)\in\{1,2\}$, with
$\Sha(E/\Q)[p^\infty]=0$, and let $L/\Qp$ be a finite tamely ramified
extension over which $E$ acquires good reduction (equivalently, over
which inertia acts trivially on the $\ell$-adic Tate module for every
prime $\ell\neq p$; reduction of type II or III is potentially good).  The hypothesis
$\GammaPD$ at $(E,p,L)$ supplies a descent map
\[
  \mathrm{descent}_p^{\mathrm{OC}}:
  LC_{\mathrm{triv}}(L_p^{\mathrm{OC}}(E_L))
  \longrightarrow
  LC_{\mathrm{triv}}(L_p^{\mathrm{OC}}(E))
\]
between trivialized leading coefficients of ordinary/anticyclotomic local
data over $L$ and over $\Qp$, a unit $\varepsilon_p(E,L)\in\Zp^\times$, and
the congruence
\[
  LC_{\mathrm{triv}}(L_p^{\mathrm{OC}}(E)) \equiv
  c_p(E)^{-1}\,\varepsilon_p(E,L)\,
  \mathrm{descent}_p^{\mathrm{OC}}
    \bigl(LC_{\mathrm{triv}}(L_p^{\mathrm{OC}}(E_L))\bigr)
  \pmod{\Zp^\times}.
\]
It also carries a factor $\mathrm{tam}_p$, which the recognition predicate
of \Cref{sec:shell-readout} matches with the Tamagawa product of $E$.
\end{definition}

Two features of this hypothesis should be noted.  First, since $p$ is odd
and $c_p(E)\in\{1,2\}$, the number $c_p(E)$ is a $p$-adic unit (in
particular $p\nmid c_p(E)$).  The congruence modulo $\Zp^\times$ therefore
holds or fails independently of whether $c_p(E)$ is $1$ or $2$: it carries
no information about $c_p(E)$.  For example, the curves $121a1$ and
$121b1$ have additive reduction at $11$ of types II and III, with $c_{11}$
equal to $1$ and $2$ respectively, and these two values have the same
class modulo $\Z_{11}^\times$.  The value of $c_p(E)$ is recorded in the
hypothesis separately.  Second, the matching of a local factor with the
global Tamagawa product is itself a comparison.  On the curve
$y^2=x^3+385x+1225$, the two places $5$ and $7$ of type III each have
$c_p=2$, and the global Tamagawa product is $4$; a single local factor is
not the global one.

\begin{definition}[$\GammaSP$: signed $\Sha$-persistence hypothesis]\label{def:closure:gamma-sha-persistence}
Let $E/\Q$, let $p$ be a prime of additive reduction at which the signed
Selmer theory used here applies, and let $L/\Qp$ be a finite tamely
ramified extension over which $E$ acquires good reduction.  Let
$\mathrm{sp}_p(E,L)$ be the signed local pairing on the cover side,
$\cores_{L/\Qp}$ the corestriction from the cover-side target to a
base-side target, and
\[
  C_{\mathrm{loc}}^{\mathrm{III}}(E,p)
  =
  \det\bigl(\cores_{L/\Qp}\circ \mathrm{sp}_p(E,L)\bigr)
\]
the determinant of the composite.  This construction, its
determinant-line convention and its applicability at $(E,p,L)$ are part
of the hypothesis; the signed theory of Kobayashi and
Pollack~\citep{Kobayashi2003,Pollack2003}, developed for good
supersingular reduction, is background for it, not a source of it.  The
hypothesis $\GammaSP$ at $(E,p,L)$ asserts
\[
  C_{\mathrm{loc}}^{\mathrm{III}}(E,p)
  \equiv c_p(E)\pmod{\Zp^\times},
\]
and carries a factor $\mathrm{sha}_p$, which the recognition predicate
matches with $\#\Sha(E/\Q)$.
\end{definition}

The same remark applies: when $p\nmid c_p(E)$, the congruence does not
see the value of $c_p(E)$.  The determinant of a composite of a map into
the cover-side target with corestriction to the base side requires a
determinant-line convention, which is part of the hypothesis, as is the
matching of $\mathrm{sha}_p$ with the global order of $\Sha$.

\begin{definition}[$\GammaGZ$: leading-term hypothesis in rank at least two]\label{def:closure:gamma-higher-gz-fixity}
Let $E/\Q$ be without complex multiplication, of analytic rank $r\ge2$.
Put
\[
  \psi_-(E,r)
  =
  \frac{L^{(r)}(E,1)}{r!}
  -
  \kapr(E)\,\Reg_{\mathrm{NT}}(E)\,\Omega_E\,\Tam(E).
\]
The hypothesis $\GammaGZ$ at $E$ asserts $\psi_-(E,r)=0$.  In the
recognition predicate it is used together with the normalization
$\kapr(E)=\#\Sha(E/\Q)/\#E(\Q)_{\tors}^{\,2}$ and with the matching of
its regulator, period and Tamagawa factors to those of $E$.
\end{definition}

The name refers to the kind of theorem that would be expected to supply
this hypothesis, a generalized Gross--Zagier formula in rank at least
$2$.  No theorem supplying this normalized higher-rank identity is
established among the inputs used here, and we retain it as a hypothesis
(see also \Cref{rmk:closure:t-e12-gap}).  It is
important to be clear about the strength of $\GammaGZ$.  Assume
$\Sha(E/\Q)$ finite, and let $\Omega_E$ and $\Tam(E)=\prod_vc_v(E)$ be as
in the BSD formula.  By the normalization of $\kapr$
\citep[Section~13]{TsiokosBSDApparatus2026}, the arithmetic term of
$\psi_-(E,r)$ is then
\[
  \kapr(E)\,\Reg_{\mathrm{NT}}(E)\,\Omega_E\,\Tam(E)
  =
  \frac{\#\Sha(E/\Q)\,\Reg_{\mathrm{NT}}(E)\,\Omega_E\,\Tam(E)}
       {\#E(\Q)_{\tors}^{\,2}},
\]
so $\GammaGZ$, read with this normalization, is the Strong BSD identity for
$E$.  The hypothesis is not weaker than the conclusion it is used for.

\begin{nonclaim}\label{nonclaim:closure:recognition-sources-not-derived}
None of the three recognition hypotheses is derived in this paper, from
the Six Birds framework or otherwise.  Each is a hypothesis about a given
curve, recorded in the formalization as a structure whose arithmetic
content is a field, not as an axiom.  We also do not claim that any of
them is necessary for the scalar conclusion: in rank at least $2$,
$\GammaGZ$ with its normalization already implies it.
\end{nonclaim}
